\documentclass[11pt,a4paper]{article}

\usepackage{geometry}
\geometry{a4paper, margin=1in}
\usepackage{xcolor}
\usepackage{colortbl}
\usepackage{booktabs}
\usepackage{tabularx}
\usepackage{enumitem}
\usepackage{setspace}
\usepackage{fancyhdr}
\usepackage{titlesec}

% Minimal fonts for XeLaTeX
\usepackage{fontspec}
\setmainfont{Calibri}

% Colors based on AIGGPA theme
\definecolor{navyblue}{RGB}{0, 51, 102}
\definecolor{gold}{RGB}{204, 153, 51}
\definecolor{sectionbg}{RGB}{240, 244, 248}

\pagestyle{fancy}
\fancyhf{}
\fancyhead[L]{\footnotesize\color{navyblue}\textbf{AIGGPA Bhopal} \quad\textbar\quad Research Rationale \& Source Explanation}
\fancyfoot[C]{\footnotesize\color{navyblue}\thepage}
\renewcommand{\headrulewidth}{0.4pt}
\renewcommand{\headrule}{\hbox to\headwidth{\color{navyblue}\leaders\hrule height \headrulewidth\hfill}}

\titleformat{\section}
  {\normalfont\Large\bfseries\color{navyblue}}
  {\thesection}{1em}{}
\titleformat{\subsection}
  {\normalfont\large\bfseries\color{navyblue}}
  {\thesubsection}{1em}{}

\setlength{\headheight}{22pt}

\begin{document}

\begin{center}
    \vspace*{1cm}
    \Huge\color{navyblue}\textbf{RESEARCH RATIONALE \&} \\[0.2em]
    \Huge\color{navyblue}\textbf{SOURCE EXPLANATION} \\[0.5em]
    \large\color{gray}Defending the Methodology, Models, and Literature \\[0.2cm]
    \normalsize\color{gold}\textbf{Prepared by:} Aryan Kori
    \vspace{1cm}
\end{center}

\section*{Introduction}
This document serves as defensively prepared material to explain the reasoning behind the methodological choices, the academic models used, the data tables presented, and the key learnings extracted from the cited literature. It is designed to provide immediate answers if the panel asks "Why did you choose this?"

\vspace{0.5cm}

\section*{1. Rationale Behind the Chosen Models}

\subsection*{The UTAUT Framework (Venkatesh et al., 2003)}
\textbf{What it is:} The Unified Theory of Acceptance and Use of Technology (UTAUT) is the world's most widely used model for understanding why individuals accept or reject new software. \\
\textbf{Why it was chosen:} Unlike older models that only looked at whether a software was useful, UTAUT considers social and organizational realities. It was chosen because it includes 'Social Influence' (do bosses push government employees to use it?) and 'Facilitating Conditions' (is there internet and technical support?). In government offices, these two factors are often more important than the software itself.

\subsection*{Thematic Analysis (Braun \& Clarke, 2006)}
\textbf{What it is:} A highly structured qualitative method for reading interview notes and sorting them into "themes" or "buckets." \\
\textbf{Why it was chosen:} For an 8-week study with 45 to 60 participants, highly complex statistical math is not viable or realistic. Thematic Analysis allows the researcher to capture the authentic, detailed frustrations of government employees and sort them logically rather than reducing human experience to a simple 1-to-5 rating scale. 

\vspace{0.5cm}

\section*{2. Reasoning Behind the Tables and Strategy}

\subsection*{The Sampling Strategy Choice}
\textbf{The Logic:} Rather than randomly surveying people, the proposal pairs Stratified Random Sampling with Purposive Sampling. \\
\textbf{The Reasoning:} Stratified sampling ensures we do not accidentally interview only junior government employees or only people in Bhopal. It forces geographical and seniority balance. Meanwhile, Purposive Sampling is used specifically for IT Officers. We must deliberately "hand-pick" IT officers because they hold the backend server knowledge that frontline government employees simply do not possess. 

\subsection*{The Estimated Budget Table (60,000 INR)}
\textbf{The Logic:} The budget is heavily weighted toward "Design and Literature Review" (32,000) and "Data Collection and Travel" (15,500). \\
\textbf{The Reasoning:} Good quality qualitative research requires physical presence. The travel budget ensures the researcher can physically view the hardware (computers/internet) on the desks of government employees, enabling triangulation (comparing what an employee says against what the researcher physically observes). The software budget (8,000) is kept minimal by relying on manual thematic coding rather than expensive enterprise analytics platforms.

\vspace{0.5cm}

\section*{3. What Was Learned from the Literature (The Sources)}

\textbf{1. The Reality of Transformation Failure (McKinsey \& Company, 2018):} \\
\textit{Learning:} Over 70 percent of digital transformations fail, and it is rarely because of bad code. It fails because of human resistance and a failure to train users. This validated the decision to study the government employees themselves rather than studying the software code.

\textbf{2. E-Government in Developing Contexts (Heeks, 2003; Ndou, 2004):} \\
\textit{Learning:} Richard Heeks established that "design-reality gaps" destroy public sector projects. When software designed in high-tech headquarters meets the reality of a rural block office with low internet, the project stalling is inevitable. This taught us we must measure 'Facilitating Conditions'.

\textbf{3. The Core Drivers of Adoption (Davis, 1989; Venkatesh, 2012):} \\
\textit{Learning:} If an employee believes a tool is too difficult to use, no amount of government mandates will make them use it effectively. They will always find workarounds (like using paper registries).

\textbf{4. National Infrastructure Baselines (DARPG, 2024; TRAI, 2024):} \\
\textit{Learning:} The Press Information Bureau and DARPG reports show that 94 percent of files in the Central Secretariat are now electronic, and internet penetration is massive. This proves the issue in Madhya Pradesh is not a lack of national capability, but a localized issue in execution and capacity building.

\vspace{0.5cm}

\section*{4. Complete Annotated Bibliography}
\textit{A brief explanation of every source cited in the proposal.}

\begin{itemize}[leftmargin=1.5em, itemsep=8pt]
    \item \textbf{Alrawabdeh (2014) | Bhatnagar (2004) | Cordella \& Bonina (2012) | Kumar \& Best (2006):} These papers provide the foundational context that public sector IT reforms must create "public value" and are much harder to sustain than regular IT projects due to bureaucracy.
    \item \textbf{Braun \& Clarke (2006):} The gold-standard academic guide on how to conduct thematic analysis properly without bias.
    \item \textbf{Davis (1989) | Venkatesh et al. (2003, 2012):} The originators of the Technology Acceptance Model and UTAUT. They provide the scientific framework for the survey questions.
    \item \textbf{Heeks (2003, 2006):} The leading global authority on why e-government fails in developing nations, emphasizing design-reality gaps.
    \item \textbf{Govt. of India (DOPT 2020 | MeitY 2015) \& PIB/DARPG Releases:} The official data sources verifying the launch of Mission Karmayogi, Digital India, and the current benchmarks for e-Office adoption.
    \item \textbf{Ndou (2004):} Highlights the unique challenges developing countries face regarding e-government, specifically physical infrastructure limits.
    \item \textbf{United Nations (2022, 2024) | World Bank (2016) | OECD (2014):} Global benchmark reports used to compare India's EGDI (E-Government Development Index) ranking against global standards.
    \item \textbf{United States GAO (2019):} Used as an international comparison point demonstrating that even advanced governments struggle with updating legacy IT systems.
\end{itemize}

\end{document}
